\chapter{Operación de modelos Akida}
\label{chap:operacion}

\section{Cargar e inspeccionar un modelo}

Los modelos Akida se distribuyen como archivos FBZ. Antes de ejecutarlos conviene conocer,
al menos, la forma de entrada, la forma de salida y el preprocesado utilizado durante su
conversión.

Carga el modelo y revisa su estructura:

\begin{lstlisting}[language=Python]
import akida

model = akida.Model("models/model.fbz")

print("input_shape:", model.input_shape)
print("output_shape:", model.output_shape)
print("ip_version:", model.ip_version)
model.summary()
\end{lstlisting}

\code{summary()} muestra las capas y la forma en la que el modelo puede ejecutarse.
La cantidad de recursos utilizados depende del propio modelo y del modo de mapeo.

\section{Seleccionar el dispositivo}

Los dispositivos físicos disponibles se obtienen mediante \code{akida.devices()}
\cite{brainchip_akida_guide}:

\begin{lstlisting}[language=Python]
import akida

devices = list(akida.devices())

if not devices:
    raise RuntimeError("No se ha detectado ningun dispositivo Akida")

device = devices[0]
print(device)
\end{lstlisting}

Un dispositivo creado con \code{akida.AKD1000()} es virtual y sirve para comprobar
el mapeo de un modelo. Para ejecutar una inferencia en la tarjeta es necesario utilizar
uno de los dispositivos devueltos por \code{akida.devices()}.

\section{Mapear el modelo en la tarjeta}

Un modelo cargado puede ejecutarse inicialmente mediante software. Para utilizar la AKD1000
debe mapearse sobre el dispositivo:

\begin{lstlisting}[language=Python]
model.map(
    device,
    hw_only=True,
    mode=akida.MapMode.AllNps
)

model.summary()
\end{lstlisting}

Con \code{hw\_only=True} se exige que la ejecución se realice completamente en el hardware
Akida \cite{brainchip_akida_guide}. Si el modelo no puede mapearse de esta forma, la llamada
produce un error.

\code{AllNps}, \code{HwPr} y \code{Minimal} son distintos modos de asignar los recursos
de la AKD1000. Para una primera ejecución puedes usar \code{AllNps}. Su efecto sobre
latencia y consumo se estudia en el \cref{chap:rendimiento}.

\section{Preparar la entrada}

La forma y el tipo de los datos dependen del modelo. Antes de ejecutar una inferencia,
comprueba que la entrada coincide con \code{model.input\_shape}.

Por ejemplo, para datos almacenados en NumPy:

\begin{lstlisting}[language=Python]
import numpy as np

data = np.load("data/inputs.npy", allow_pickle=False)

print("shape:", data.shape)
print("dtype:", data.dtype)
print("min/max:", data.min(), data.max())

if tuple(data.shape[1:]) != tuple(model.input_shape):
    raise ValueError(
        f"Forma de entrada {data.shape[1:]} "
        f"!= {model.input_shape}"
    )
\end{lstlisting}

El tipo, el rango y el orden de los canales deben ser los mismos utilizados durante
la preparación y conversión del modelo. No existe un rango de entrada común para todos
los modelos Akida.


\section{Ejecutar el modelo}

La API de Akida permite ejecutar el modelo mediante \code{forward()} o \code{predict()}.
La elección depende de cómo se haya definido la salida del modelo y de cómo se haya
validado durante su conversión \cite{brainchip_akida_guide}.

Por ejemplo:

\begin{lstlisting}[language=Python]
output = model.forward(data)
print(output)
\end{lstlisting}

Durante una campaña de medidas debe mantenerse el mismo método de ejecución para que
los resultados sean comparables.

También puedes usar la herramienta incluida en este repositorio:

\begin{lstlisting}[language=bash]
akd1000-bringup run \
  --model models/model.fbz \
  --input data/input.npy \
  --method forward \
  --clock Performance \
  --map AllNps \
  --output outputs/raw-output.npy
\end{lstlisting}

La herramienta carga la entrada, ejecuta el modelo y guarda la salida en un archivo NumPy.
La entrada debe tener la forma

\begin{lstlisting}
(batch, *model.input_shape)
\end{lstlisting}

Los modelos con varias entradas o salidas requieren tratarlos directamente mediante
la API de Akida.

\section{Modo de reloj}

La AKD1000 dispone de los modos de reloj \code{Performance}, \code{Economy} y
\code{LowPower}. El modo puede seleccionarse mediante la API:

\begin{lstlisting}[language=Python]
device.soc.clock_mode = akida.soc.ClockMode.Performance
print(device.soc.clock_mode)
\end{lstlisting}

El modo de reloj debe anotarse cuando se comparen medidas de rendimiento. Un modo de menor
potencia instantánea no implica necesariamente una menor energía por inferencia, ya que también
puede cambiar el tiempo de ejecución.